% !TeX root = ../technical_paper.tex

\section{总结与展望}

\subsection{工作总结}

本文围绕面向光储融合应用的多MCU协同智能电源管理平台开展研究，完成了控制平台搭建、BMS系统联调验证、无线监控实现和三端口功率级硬件设计。系统采用GD32G553主控、GD32C113 BMS和GD32VW553无线模块分层协同架构，实现了控制、保护、通信和监控的解耦。

实验结果表明：电压采样最大误差$+4\,\mathrm{mV}$，电流采样$5\,\mathrm{A}$以内误差控制在$50\,\mathrm{mA}$，SOC估算最大偏差约$4.2\%$，CAN通信丢包率约$0.067\%$，均衡功能将单体电压差从$100\,\mathrm{mV}$降至$20\,\mathrm{mV}$。CC/CV充电切换正常，系统能够根据BMS安全状态正确切换工作模式，上位机和App可完成状态显示与远程控制。

三端口功率变换平台已完成硬件设计与PV$\rightarrow$Battery基础功能验证，但PV$\rightarrow$Load、Battery$\rightarrow$Load及多模式能量流控制仍需后续阶段完善。现有工作为后续闭环控制和整机联调奠定了软硬件基础。

\subsection{后续工作}

后续将在现有控制、采样、保护和通信平台基础上，逐步开展以下工作：

\begin{enumerate}[label=\arabic*),leftmargin=*,itemsep=0.3ex]
  \item 完成三端口功率级实物闭环测试，验证额定工况下的效率、温升和动态响应
  \item 增加不同SOC、负载和温度条件下的测试样本，完善系统性能评价数据
  \item 优化无线监控界面，增加历史数据记录、故障日志导出和曲线显示功能
  \item 建立更完整的故障注入测试流程，提高系统在复杂工况下的安全性和可靠性
\end{enumerate}

现有成果已为后续主功率级接入、闭环控制调试和整机性能测试提供了可靠基础。
